%% Đề 40 -- de-tuyen-sinh-lop-10-thpt-mon-toan-nam-2012-2013-so-gddt-lang-son.pdf
\deheader{ĐỀ SỐ 40}%
  {Kỳ thi Tuyển sinh vào lớp 10 THPT — Năm học 2012--2013}%
  {Sở GD\&ĐT Lạng Sơn | Chỉ phần Đại số}
\addcontentsline{toc}{section}{Đề 40 — Tuyển sinh vào lớp 10 THPT — Năm học 2012--2013}

\begin{de_cau}{Câu 1}{2,0 điểm}
\begin{enumerate}[label=\alph*)]
  \item Rút gọn biểu thức: $A = \sqrt{12} - \sqrt{27} + \sqrt{48}$.
  \item Giải phương trình: $3x - 6 = 0$.
\end{enumerate}
\end{de_cau}

\begin{loi_giai_cau}
\textbf{a)} $A = 2\sqrt{3} - 3\sqrt{3} + 4\sqrt{3} = \boxed{3\sqrt{3}}$.

\textbf{b)} $3x=6\Rightarrow\boxed{x=2}$.
\end{loi_giai_cau}

\begin{de_cau}{Câu 2}{2,0 điểm}
Cho biểu thức:
\[P = \left(\dfrac{\sqrt{x}}{\sqrt{x}-1} - \dfrac{1}{x-\sqrt{x}}\right) : \dfrac{\sqrt{x}+1}{3}\]
với $x > 0,\; x \neq 1$.
\begin{enumerate}[label=\alph*)]
  \item Rút gọn $P$.
  \item Tìm $x$ để $P = 1$.
\end{enumerate}
\end{de_cau}

\begin{loi_giai_cau}
\textbf{a)} $\dfrac{\sqrt{x}}{\sqrt{x}-1}-\dfrac{1}{\sqrt{x}(\sqrt{x}-1)}=\dfrac{x-1}{\sqrt{x}(\sqrt{x}-1)}=\dfrac{(\sqrt{x}-1)(\sqrt{x}+1)}{\sqrt{x}(\sqrt{x}-1)}=\dfrac{\sqrt{x}+1}{\sqrt{x}}$.

$P=\dfrac{\sqrt{x}+1}{\sqrt{x}}\cdot\dfrac{3}{\sqrt{x}+1}=\boxed{\dfrac{3}{\sqrt{x}}}$.

\textbf{b)} $P=1\Rightarrow\dfrac{3}{\sqrt{x}}=1\Rightarrow\sqrt{x}=3\Rightarrow\boxed{x=9}$ (thỏa mãn).
\end{loi_giai_cau}

\begin{de_cau}{Câu 3}{2,0 điểm}
Cho phương trình $x^2 - 2(m+2)x + m^2 + 4 = 0$.
\begin{enumerate}[label=\alph*)]
  \item Giải phương trình khi $m = 1$.
  \item Tìm $m$ để phương trình có hai nghiệm $x_1, x_2$ thỏa mãn $x_1^2 + x_2^2 = 16$.
\end{enumerate}
\end{de_cau}

\begin{loi_giai_cau}
\textbf{a)} $m=1$: $x^2-6x+5=0\Rightarrow(x-1)(x-5)=0\Rightarrow\boxed{x_1=1,\; x_2=5}$.

\textbf{b)} $\Delta'=(m+2)^2-(m^2+4)=4m\ge0\Rightarrow m\ge0$.

Vi-ét: $x_1+x_2=2(m+2)$, $x_1x_2=m^2+4$.

$x_1^2+x_2^2=(x_1+x_2)^2-2x_1x_2=4(m+2)^2-2(m^2+4)=2m^2+16m+8=16$.

$2m^2+16m-8=0\Rightarrow m^2+8m-4=0\Rightarrow m=-4\pm2\sqrt{5}$.

Kết hợp $m\ge0$: $m=\boxed{-4+2\sqrt{5}}$.
\end{loi_giai_cau}

\begin{de_cau}{Câu 5}{0,5 điểm}
Cho $a, b$ là các số thực thỏa mãn $a > 0$, $b > 0$ và $a + b \le 2$. Tìm giá trị nhỏ nhất của biểu thức:
\[M = \left(a+\dfrac{1}{a}\right)^2 + \left(b+\dfrac{1}{b}\right)^2.\]
\end{de_cau}

\begin{loi_giai_cau}
Hàm $f(t)=\left(t+\dfrac{1}{t}\right)^2$ lồi và giảm trên $(0,1]$. Với $a+b\le2$ và $a,b>0$, cực tiểu đạt khi $a=b=1$:

$M_{\min}=\left(1+1\right)^2+\left(1+1\right)^2=4+4=\boxed{8}$.
\end{loi_giai_cau}
